\lesson{6}{Sep 27 2021 Mon (07:30:12)}{Polynomial Operations}{Unit 1}

\subsubsection{Adding Polynomials}
\label{sub_sub_sec:adding_polynomials}

\begin{definition}[Like Terms]
  \label{def:like_terms}

  Like terms are terms that have the same \textbf{variable} and
  \textbf{exponent}. 

  Ex: $5y^{2}$ and $5909y^{2}$. \\
  Ex: $9zsa^{d}$ and $8sa^{d}z$.
\end{definition}

The most important part of adding or subtracting polynomials is identifying \textbf{like
terms} (\ref{def:like_terms}).

\begin{example}
  \label{exm:adding_polynomials}
  Let's add the following polynomials: $(7x^{2} + 3x - 5) + (2x^{2} - 4x + 9)$

  \begin{equation}
    \label{eqt:adding_polynomials}
    \begin{split}
      \textrm{Step $1$:} \quad &\textrm{Distribute any Coefficient} \\
      1(7x^{2} + 3x - 5) + 1(2x^{2} - 4x + 9) &= 7x^{2} + 3x - 5 + 2x^{2} - 4x + 9 \\
      \textrm{Step $2$:} \quad &\textrm{Identify the Like Terms} \\
      7x^{2} + 3x - 5 + 2x^{2} - 4x + 9 &= \boxed{9x^{2} - x + 4} \\
    \end{split}
  .\end{equation}

  The reason there's a $1$ in front of each polynomial is because we're adding,
  and when you add, you distribute a positive $1$. Take a look at this
  polynomial: $5(3x^{2} + 3x - 1) + 4(x^{2} + 2)$. You distribute a positive $5$ 
  to the first polynomial and a positive $4$ to the second polynomial.
\end{example}

% subsubsection adding_polynomials (end)

\subsubsection{Subtracting Polynomials}
\label{sub_sub_sec:subtracting_polynomials}

\begin{review}[Distributive Property]
  \label{rev:distributive_property}

  \begin{equation}
    \label{eqt:distributive_property}
    \begin{split}
      a(b + c) &= a(b) + a(c) \\
      -a(b + c) &= -(a(b)) + -(a(c)) \\
    \end{split}
  .\end{equation}
\end{review}

Subtracting polynomials is almost exactly like adding polynomials. The only
difference is you distribute a negative coefficient instead of a positive one.

\begin{example}
  \label{exm:subtracting_polynomials}
  Let's subtract the following polynomials:
  $(9x^{2} + x + 2) - (7x^{2} + 3x - 4)$.

  \begin{equation}
    \label{eqt:subtracting_polynomials}
    \begin{split}
      \textrm{Step $1$:} \quad &\textrm{Distribute any Coefficient} \\
      (9x^{2} + x + 2) - (7x^{2} + 3x - 4) &= 9x^{2} + x + 2 - 7x^{2} - 3x + 4 \\
      \textrm{Step $2$:} \quad &\textrm{Identify the Like Terms} \\
      9x^{2} + x + 2 - 7x^{2} - 3x + 4 &= \boxed{2x^{2} - 2x + 6}
    \end{split}
  .\end{equation}

  This time, instead of distributing a positive $1$, we distribute a negative
  $1$ because we're subtracting. Take a look at the \textbf{Distributive
  Property} review (\ref{rev:distributive_property}) if you need a refresher.
\end{example}

% sub subsection subtracting_polynomials (end)

\subsubsection{Using the Distributive Property}
\label{sub_sub_sec:using_the_distributive_property}

\paragraph{Multiplying by a Monomial}
\label{par:multiplying_by_a_monomial}

\begin{definition}[Monomial]
  \label{def:monomial}

  A monomial is a polynomial with only one term. Ex:
  $\quad 7 \quad 4x \quad 19z^{10}$.
\end{definition}

\begin{example}
  \label{exm:multiplying_by_a_monomial}

  Let's try to multiply the following monomial by the following polynomial:
  $3x^{2}(5x^{2} - 7x + 8)$
  
  \begin{equation}
    \label{eqt:multiplying_by_a_monomial}
    \begin{split}
      3x^{2}(5x^{2} - 7x + 8) &= 3x^{2} \times 5x^{2} \\
                        &= 15x^{4} \\
      3x^{2}(5x^{2} - 7x + 8) &= 3x^{2} \times (-7x) \\
                        &= -21x^{3} \\
      3x^{2}(5x^{2} - 7x + 8) &= 3x^{2} \times 8 \\
                        &= 24x^{2} \\
                        &= \boxed{15x^{4} - 21x^{3} + 24x^{2}} \\
    \end{split}
  .\end{equation}
\end{example}

% paragraph multiplying_by_a_monomial (end)

\paragraph{Multiplying by a Binomial}
\label{par:multiplying_by_a_binomial}

\begin{definition}[Binomial]
  \label{def:monomial}

  A binomial is a polynomial with only two term. Ex:
  $\quad 7 + 9x \quad 4y^{3} - 4x \quad 19z^{10} \times 8z$.
\end{definition}

\begin{example}
  \label{exm:multiplying_by_a_binomial}

  Let's try to multiply the following monomial by the following polynomial:
  $(x - 7)^{2}$. The following polynomial can also be written as: $(x - 7)(x - 7)$
  
  \begin{equation}
    \label{eqt:multiplying_by_a_binomial}
    \begin{split}
      (x - 7)(x - 7) &= x \times x \\
                        &= x^{2} \\
      (x - 7)(x - 7) &= x \times (-7) \\
                        &= -7x \\
      (x - 7)(x - 7) &= (-7) \times x \\
                        &= -7x \\
      (x - 7)(x - 7) &= (-7) \times (-7) \\
                        &= 49 = x^{2} - 7x - 7x + 49 \\
                        &= \boxed{x^{2} - 14x + 49}
    \end{split}
  .\end{equation}
\end{example}

% paragraph multiplying_by_a_binomial (end)

\paragraph{Multiplying by a Trinomial}
\label{par:multiplying_by_a_trinomial}

\begin{definition}[Trinomial]
  \label{def:trinomial}

  A trinomial is a polynomial with three terms. Ex:
  $\quad 8x + 1 - y \quad u + 8x - 1 \quad 1 \times 14x - z$
\end{definition}

\begin{example}
  \label{exm:multiplying_by_a_trinomial}

  Let's try to multiply the following polynomial by the following binomial
  $(4x^{2} - 5x + 9)(x - 5)$
  
  \begin{equation}
    \label{eqt:multiplying_by_a_trinomial}
    \begin{split}
      (4x^{2} - 5x + 9)(x - 5) &= 4x^{2} \times x \\
                            &= 4x^{3} \\
      (4x^{2} - 5x + 9)(x - 5) &= 4x^{2} \times (-5) \\
                            &= -20x^{2} \\
      (4x^{2} - 5x + 9)(x - 5) &= (-5x) \times x \\
                            &= -5x^{2} \\
      (4x^{2} - 5x + 9)(x - 5) &= (-5x) \times (-5) \\
                            &= 25x \\
      (4x^{2} - 5x + 9)(x - 5) &= 9 \times x \\
                            &= 9x \\
      (4x^{2} - 5x + 9)(x - 5) &= 9 \times (-5) \\
                            &= -45 = 4x^{3} - 20x^{2} - 5x + 25x + 9x - 45 \\
                            &= \boxed{4x^{3} - 25x^{2} + 34x - 45}
    \end{split}
  .\end{equation}
\end{example}

% paragraph multiplying_by_a_trinomial (end)

% subsubsection using_the_distributive_property (end)

\subsubsection{Equations vs Functions}
\label{sub_sub_sec:equations_vs_functions}

\begin{definition}[Equation]
  \label{def:equation}

  An equation is usually defined as the state of being equal and is often shown
  as a math expression with equal values on either side, or refers to a problem
  where many things need to be taken into account. Ex: $2 + 2 = 4 + 2 - 1(2)$
\end{definition}

\begin{definition}[Function]
  \label{def:function}

  A function in mathematics is an expression, rule, or law that defines a
  relationship between one variable and another variable. Ex: $f(x) = x^{2} + x$

  In order for an equation to be a function, there must be only one $y$ value
  for each $x$ value. The function $f(x) = x^{2}$ isn't a function. If you input
  $f(2) = 2^{2} = 4$ and $f(-2) = -2^{2} = 4$. You get the same $y$ value for
  two different $x$ values. So, it doesn't meet the requirements to be a
  function.
\end{definition}

The mathematical term function (\ref{def:function}) applies to some equations
(\ref{def:equation}) in math.
But, not all equations are functions.

\begin{example}
  \label{exm:equations_vs_functions_1}

  $y = 2x + 3$ 

  Let's take a look at the equations $x/y$ table:

  \begin{center}
    \begin{tabular}{|r|l|}
      \hline
      x & y \\ \hline
      -2 & -1 \\ \hline
      -1 & 1 \\ \hline
      0 & 3 \\ \hline
      1 & 5 \\ \hline
      2 & 7 \\ \hline
    \end{tabular}
  \end{center}

  Since this equation doesn't have any $y$ values for multiple $x$ values, then
  it's considered a function.
\end{example}

% subsubsection equations_vs_functions (end)

\subsubsection{Vertical Line Test}
\label{sub_sub_sec:vertical_line_test}

\begin{definition}[Vertical Line Test]
  \label{def:vertical_line_test}

  The \textbf{Vertical Line Test} is a visual method to determine if an equation
  is a function or not. You first must graph the equation. Then, if you two or
  more $y$ intercepts, then it's not a function because it has the same $y$ 
  value for multiple $x$ values.

  Likewise, if it has $1$ or $0$ $y$ intercepts, then it's a function.
\end{definition}

Another way to determine of an equation is a function is to graph the equation
and perform the \textbf{Vertical Line Test} (\ref{def:vertical_line_test}).

\begin{example}
  \label{exm:vertical_line_test}

  Let's graph the equation: $y = 2x + 3$:

  \begin{figure}[H]
    \centering
    \incfig{vertical_line_test}
    \caption{Vertical Line Test}
    \label{fig:vertical_line_test}
  \end{figure}

  Since there aren't any multiple $y$ intercepts, then it's a function.
\end{example}

% subsubsection vertical_line_test (end)

\subsubsection{Inverses of Functions}
\label{sub_sub_sec:inverses_of_functions}

Finding the inverse of an integer means performing an operating that will cancel
that number. For example, if you are given the number $7$, you could add $-7$ to
it to cancel it: $7 + (-7) = 0$. Or you could multiply it by it's reciprocal:
$7 \times \frac{1}{7} = 1$.

With equations and functions, finding the inverse requires reversing the
variables and solving for $y$.

\begin{example}
  \label{exm:inverses_of_functions}

  Let $f(x) = 3x - 5$. Find its inverse represented by the function notation:
  $f^{-1}(x)$ 

  \begin{equation}
    \label{eqt:inverses_of_functions}
    \begin{split}
      f(x) &= 3x - 5 \\
      y &= 3x - 5 \\
      x &= 3y - 5 \\
      x + 5 &= 3y \\
      \frac{x + 5}{3} &= \frac{3y}{3} \\
      \frac{x + 5}{3} &= y \\
      f^{-1}(x) &= \frac{x + 5}{3} \\
    \end{split}
  .\end{equation}
\end{example}

% subsubsection inverses_of_functions (end)

\subsubsection{Operations on Functions}
\label{sub_sub_sec:operations_on_functions}

\paragraph{Addition}
\label{par:addition}

The addition of two functions $f(x)$ and $g(x)$ is represented using the
following notation: $f(x) + g(x)$

\begin{example}
  \label{exm:addition}

  Let $f(x) = 4x - 7$
  
  Let $g(x) = 10x - 3$
  
  \begin{equation}
    \label{eqt:addition_1}
    \begin{split}
      f(x) &= 4x - 7 \\
      g(x) &= 10x - 3 \\
      f(x) + g(x) &= (4x - 7) + (10x - 3) \\
                  &= 4x - 7 + 10x - 3 \\
                  &= \boxed{14x - 10} \\
    \end{split}
  .\end{equation}
\end{example}

% paragraph addition (end)

\paragraph{Subtraction}
\label{par:subtraction}

Subtracting functions is similar to adding functions. The only difference is
that you need to be very caucus with the signs.

\begin{example}
  \label{exm:subtraction}

  Let $f(x) = 4x - 7$
  
  Let $g(x) = 10x - 3$
  
  \begin{equation}
    \label{eqt:addition_1}
    \begin{split}
      f(x) &= 4x - 7 \\
      g(x) &= 10x - 3 \\
      f(x) - g(x) &= (4x - 7) - (10x - 3) \\
                  &= 4x - 7 - 10x + 3 \\
                  &= \boxed{-6x - 4} \\
    \end{split}
  .\end{equation}
\end{example}

% paragraph subtraction (end)

\paragraph{Multiplication}
\label{par:multiplication}

The multiplication of functions is usually just an extension of the distribute
property (\ref{rev:distributive_property}). The notation used to multiply
functions is $f(x) \times g(x)$.

\begin{example}
  \label{exm:multiplication}

  Let $f(x) = 4x - 7$

  Let $g(x) = 10x - 3$ 

  \begin{equation}
    \label{eqt:multiplication}
    \begin{split}
      f(x) \times g(x) &= (4x - 7)(10x - 3) \\
                  &= 4x \times 10x = 40x^{2} \\
                  &= 4x \times (-3) = -12x \\
                  &= (-7) \times 10x = -70x \\
                  &= (-7) \times (-3) = 21 \\
                  &= 40x^{2} - 12x - 70x + 21 \\
                  &= \boxed{40x^{2} - 82x + 21} \\
    \end{split}
  .\end{equation}

\end{example}

% paragraph multiplication (end)

% subsubsection operations_on_functions (end)

\subsubsection{Function Composition of Polynomials}
\label{sub_sub_sec:function_composition_of_polynomials}

\begin{definition}[Function Composition]
  \label{def:function_composition}

  \textbf{Function Composition} is applying the function to the result of
  another.  

  There are multiple ways for writing function compositions. Here's one way:
  $(f \circ g)(x)$. This is the same thing as: $f(g(x))$. 

  What this means is we solve for $g(x)$, then whatever the output of that is, we
  plug that into $f(x)$.
\end{definition}

\begin{example}
  \label{exm:function_composition_of_polynomials_1}

  Given the functions $f(x) = 2x^{2} - 4x + 1$ and $g(x) = x^{2} + 3$, find
  $(f \circ g)(x)$.
  
  \begin{equation}
    \label{eqt:function_composition_of_polynomials}
    \begin{split}
      \textrm{Step 1:} \quad & \quad \textrm{Substitute expression} \\
      f(x) = 2x^{2} - 4x + 1 \quad & \quad g(x) = x^{2} + 3 \\
      f(g(x)) &= 2(x^{2} + 3)^{2} - 4(x^{2} + 3) + 1 \\
      \textrm{Step 2:} \quad & \quad \textrm{Simplify} \\
                & \textrm{Rewrite $(x^{2} + 3)^{2}$ as $(x^{2} + 3)(x^{2} + 3)$} \\
              &= 2(x^{2} + 3)(x^{2} + 3) - 4(x^{2} + 3) + 1 \\
                & \textrm{Multiply $(x^{2} + 3)(x^{2} + 3)$} \\
              &= 2(x^{4} + 3x^{2} + 3x^{2} + 9) - 4(x^{2} + 3) + 1 \\
              &= 2(x^{4} + 6x^{2} + 9) - 4(x^{2} + 3) + 1 \\
                & \textrm{Distribute the $2$ and $-4$} \\
              &= 2x^{4} + 12x^{2} + 18 - 4x^{2} - 12 + 1 \\
                & \textrm{Combine like terms} \\
              &= \boxed{2x^{4} + 8x^{2} + 7} \\
    \end{split}
  .\end{equation}
\end{example}

Function composition (\ref{def:function_composition}) is different from function
multiplication. By substituting the function $g(x)$ into the function $f(x)$,
you are evaluating $f(x)$ at $g(x)$. As $g(x)$ changes for each $x$ value, the
composition of $f(g(x))$ will also change because you will be evaluating $f(x)$
at the new $g(x)$ value each time.

If you're given a specific value for $g(x)$, then you could simply substitute
that value into $f(x)$ to solve for the composition, $f(x)$.

\begin{example}
  \label{exm:function_composition_of_polynomials_2}

  If you are given $g(4) = 10$ and you're asked to find $f(g(4))$, you would
  substitute $10$ into $f(x)$. You're evaluating $f(x)$ at the answer for
  $g(4)$. Most often, you will need to write a new function, $f(g(x))$, using
  the original functions so that you can have a new function that will represent
  all $x$ values.
\end{example}

Paying attention to the units that each function, $f(x)$ and $g(x)$, is measured
in will help you figure out what represent. If $f(x)$ is measured in
$\frac{\textrm{dollars}}{\textrm{bagel}}$ and $g(x)$ is measured in
$\frac{\textrm{bagels}}{\textrm{hour}}$, then the composition of $f(g(x))$ is taking a function
measured in $\frac{\textrm{bagels}}{\textrm{hour}}$ and substituting that into a
function being measured in $\frac{\textrm{dollars}}{\textrm{bagel}}$.
This means the composition $f(g(x))$ will be measured in
$\frac{\textrm{dollars}}{\textrm{hour}}$

% subsubsection function_composition_of_polynomials (end)

\subsubsection{Verifying Inverses}
\label{sub_sub_sec:verifying_inverses}

Function composition is often used to verify inverse functions. If $f(g(x)) =
g(f(x)) = x$, then the functions are inverses of each other.

\begin{example}
  \label{exm:verifying_inverses}

  Verify that $f(x) = \frac{x + 5}{3}$ is the inverse of $g(x) = 3x - 5$:

  Use function composition to evaluate $(f \circ g)(x)$:

  \begin{equation}
    \label{eqt:verifying_inverses}
    \begin{alignedat}{3}
      (f \circ g)(x) &= \frac{(3x - 5) + 5}{3} && \\
                 &= \frac{3x}{3} = \boxed{x} && \\
      \textrm{Step 1:} \quad  & \quad \textrm{Now Simplify} && \\
                 &= 3(\frac{x + 5}{3}) - 5  &&= \frac{3(x) + 3(5)}{3} - 5 \\
                                          & &&= \frac{3x + 15}{3} - 5 \\
                                          & &&= \frac{3x}{3} + 5 - 5 \\
                                          & &&= \frac{3x}{3} \\
                                          & &&= \boxed{x} \\
    \end{alignedat}
  \end{equation}

  Because $(f \circ g)(x) = (g \circ f)(x) = x$, then $f(x)$ and $g(x)$ are
  inverses.
\end{example}

% subsubsection verifying_inverses (end)

\newpage

